%% 自動生成（gen.py）
\bktocchap{4-0}{ロボットの手は、どこにある？}{長さと角度がわかれば、場所が決まる}{4}{2}{2}{8}{J}
\bktocpart{第1部}{図形}{24}
\bktocchap{4-1}{きまりから考えを積み上げる：定義・公理・定理}{三角形の角をたすと、なぜ必ず180°？}{6}{2}{3}{26}{C}
\bktocchap{4-2}{角度と円、ラジアン}{角度を「長さ」ではかる方法}{4}{1}{2}{56}{C}
\bktocchap{4-3}{にた形と比：相似}{影の長さから、木の高さがわかる}{6}{2}{3}{98}{C}
\bktocchap{4-4}{三平方の定理}{ななめの距離を、たてと横から計算する}{6}{3}{3}{128}{C}
\bktocchap{4-5}{面積と体積：細かく分けてたす}{円を切って並べかえると、長方形になる？}{2}{1}{2}{160}{C}
\bktocpart{第2部}{三角比と三角関数}{194}
\bktocchap{4-6}{坂道と三角比}{坂の上の荷物にかかる力を、2つに分ける}{4}{2}{3}{196}{C}
\bktocchap{4-7}{三角形を解く：正弦定理と余弦定理}{行けない場所までの距離をはかる}{4}{3}{3}{234}{C}
\bktocchap{4-8}{三角関数：円運動の影}{回る点の影は、行ったり来たりする}{4}{3}{4}{268}{C}
\bktocchap{4-9}{波をたし合わせる：加法定理とうなり}{2つの音を重ねると、ウワンウワンと聞こえる}{4}{4}{4}{306}{C}
\bktocpart{第3部}{ベクトル}{338}
\bktocchap{4-10}{向きと大きさをもつ量：ベクトル}{2人で引くと、荷物はどっちへ動く？}{6}{2}{2}{340}{C}
\bktocchap{4-11}{ベクトルを成分で計算しよう}{川に流されながら渡る船は、どこに着く？}{6}{2}{3}{374}{C}
\bktocchap{4-12}{内積：どれだけ同じ向きか}{進む向きと直角に押しても、仕事は0}{6}{3}{4}{408}{C}
\bktocchap{4-13}{空間の座標と3次元のベクトル}{部屋の中の1点を、3つの数で表す}{6}{4}{4}{434}{C}
\bktocchap{4-14}{外積：回す力と右ねじ}{ドアのノブは、なぜちょうつがいから遠い？}{6}{4}{5}{468}{C}
\bktocchap{4-15}{いろいろな座標：極座標・円柱座標・球座標}{地球の上の場所は、緯度と経度で表す}{6}{4}{5}{500}{C}
\par\vspace{2mm}
\bktocbacklist{\bkbi{こたえ}{526}\bkbi{\R{記号}{き|ごう}カード}{539}\bkbi{よく\R{使}{つか}う\R{定数}{てい|すう}・\R{単位}{たん|い}・\R{公式}{こう|しき}}{543}\bkbi{できるようになったことチェック\R{表}{ひょう}}{550}\bkbi{つぎの\R{巻}{かん}へ}{553}\bkbi{さくいん}{556}}
